\section{A PTAS for uniform shallow caps in convex position}
\label{sec:uniform-cap-ptas}

The bounded-cap restriction also permits an approximation scheme for
uniform expected depth when the hypotheses are in convex position.
The construction first clears short gaps in their cyclic order.
Every remaining query then acts inside one bounded-size block.
A count-vector dynamic program retains all choices of how to interleave
queries from these blocks. Uniform weights make blocks with the same
local test family interchangeable.

\begin{theorem}\label{thm:uniform-cap-ptas}
Suppose $\HH=P\cup\{p\}$ is in strictly convex position, $p$ is a
distinguished hypothesis, and every supplied side missing $p$ contains
at most a fixed number $b\ge1$ of hypotheses. Under the uniform prior,
expected depth admits a polynomial-time approximation scheme. Every
query in the returned tree is supplied; singleton caps need not be
available.

More explicitly, for $0<\varepsilon\le1$, put
\[
 \ell=\max\left\{2b,\ b(b-1),\
                 \left\lceil\frac{8b^2}{\varepsilon}\right\rceil\right\},
 \qquad d_{b,\ell}=2^{(b+1)\ell}.
\]
One running-time bound is
$f(b,\ell)(n+1)^{d_{b,\ell}}\operatorname{poly}(n,m,L)$, where
$L$ is the input bit length and $f$ depends only on $b,\ell$.
The polynomial exponent depends on $b$ and $\varepsilon$; the
guarantee is a PTAS.
\end{theorem}

\begin{proof}
Write $N=|P|=n-1$ and work with the raw uniform cost
$G(T)=\sum_{h\in\HH}d_T(h)$; dividing by $n$ preserves ratios.
Write $\OPT_G=\min_TG(T)=n\OPT_E$.
When $N<\ell$, exact subset Bellman optimization is sufficient,
since the state size depends only on $b$ and $\varepsilon$.
Assume henceforth that $N\ge\ell$.

\paragraph{Clear gaps between small blocks.}
Each cap is a cyclic interval in the convex order of $P$: a line
crosses the boundary of its convex polygon at most twice. Partition
that order into consecutive pieces of length $\ell$, with a possibly
shorter last piece. Mark the first $b$ points of each piece, or the
whole piece when its size is at most $b$. Let $D$ be the marked set
and $r=|D|$. Then
\begin{equation}\label{eq:uniform-cap-deletion-count}
 r\le b\lceil N/\ell\rceil\le\frac{2bN}{\ell}.
\end{equation}
The unmarked portions are blocks of at most $\ell$ points. Between
neighboring nonempty blocks, including across the cyclic wraparound,
there are at least $b$ marked points. A cyclic cap of size at most $b$
cannot meet two such blocks, because it would also contain an entire
intervening gap and hence at least $b+2$ points.

While a marked point $x$ remains on the miss spine, query any supplied
cap containing $x$. Such a cap exists because the supplied tests
distinguish $x$ from $p$. Identify its first-hit group by a proper
local supplied-query tree, and continue on the miss side. Separation
gives a local height at most $|A|-1\le b-1$ for each hit group $A$:
repeatedly query a line separating two remaining hypotheses. Each
spine query clears a new marked point, so at most $r$ are used.
Additional unmarked points may leave through these exits. The
residual portions of the blocks still have size at most $\ell$, and
every residual cap lies in one block.

The prefix's spine nodes cost at most $r(N+1)$ in total. At most
$br$ points leave through its hit exits, each with at most $b-1$
additional local queries. Thus the complete prefix contribution is
\begin{equation}\label{eq:uniform-cap-prefix-cost}
 G_{\rm prefix}\le r(N+1)+b(b-1)r\le2r(N+1),
\end{equation}
using $N\ge\ell\ge b(b-1)$.

\paragraph{Optimize all block interleavings exactly.}
Relabel the points of each residual block by their positions. Its
type $T$ is its size together with the family $\mathcal F_T$ of all
distinct nonempty supplied-query masks on that block. No cap can
wrap around a block's ends, since it would contain an entire marked
gap. Deleting additional points preserves interval order. Thus every
mask is an ordinary interval of at most $b$ positions. Padding to
$\ell$ positions gives at most $b\ell$ possible masks, hence at most
$2^{b\ell}$ families. The union of a family's masks is its full block,
since every point is separated from $p$, so the family also determines
the block size. A local
state $(T,R)$ records a nonempty live mask $R$ in a block of type $T$.
There are at most $d_{b,\ell}$ local states. Let $c_{T,R}$ count the
blocks in each state, and let
\[
 W(c)=1+\sum_{T,R}c_{T,R}|R|
\]
be the total live mass, including $p$. Uniformity ensures that these
counts retain all information needed to optimize the continuation.

For each type, precompute by Algorithm~\ref{alg:exact-tree} the exact
raw local cost $g_T(A)$ of identifying every subset $A$ using its
supplied masks. Here the positive branch has no distinguished
witness, and $g_T(A)=0$ for $|A|\le1$. Let $e_{T,R}$ be a unit count
vector, with $e_{T,\varnothing}=0$. The exact miss-continuation
recurrence is
\begin{equation}\label{eq:uniform-cap-count-dp}
 \begin{aligned}
 F(0)&=0,\\
 F(c)&=W(c)+\min_{\substack{c_{T,R}>0,\ q\in\mathcal F_T\\
                          A=R\cap q\ne\varnothing}}
 \left\{g_T(A)+F\bigl(c-e_{T,R}+e_{T,R\setminus A}\bigr)\right\}.
 \end{aligned}
\end{equation}
Every transition removes a live point, so the recurrence is acyclic.
At any proper query, its hit group lies in one block and its miss
branch retains $p$ and the other blocks. The positive subtree can be
optimized locally, giving precisely the recurrence's lower bound
for every adaptive tree. Conversely, a transition chooses a physical
block of the specified type and live mask, and an original supplied
query realizing $q$. That query misses every other residual block.
Its stored local positive code and miss continuation realize the
recurrence. Backtracking therefore reconstructs a legal identifying
tree, including all choices of block interleaving.

Each count lies between zero and $N$, giving at most
$(N+1)^{d_{b,\ell}}$ states. Local Bellman tables and the transitions per
state depend only on $\ell$, and constructing the types takes
polynomial time in the input. All raw costs are integral and at most
$n(n-1)$, so their bit complexity is polynomial.

\paragraph{Charge the prefix against the optimum.}
Restrict an optimal full tree to the residual hypotheses, suppressing
constant-answer nodes. Restriction can only decrease depths, so the
exact residual cost from \eqref{eq:uniform-cap-count-dp} is at most
$\OPT_G$. On the other hand, the witness spine of any full tree has
first-hit groups of size at most $b$, and hence
\begin{equation}\label{eq:uniform-cap-rank-lower}
 \OPT_G\ge\sum_{i=1}^N\lceil i/b\rceil
          \ge\frac{N(N+1)}{2b}.
\end{equation}
Combining \eqref{eq:uniform-cap-prefix-cost},
\eqref{eq:uniform-cap-deletion-count}, and
\eqref{eq:uniform-cap-rank-lower} gives
\[
 \frac{G_{\rm alg}-\OPT_G}{\OPT_G}
 \le\frac{4br}{N}
 \le\frac{8b^2}{\ell}
 \le\varepsilon.
\]
Dividing by $n$ proves the expected-depth guarantee.
\end{proof}

\begin{algorithm}[H]
\caption{Uniform expected-depth PTAS for bounded convex-position caps}
\label{alg:uniform-cap-ptas}
\KwIn{Uniform hypotheses $P\cup\{p\}$ in convex position; supplied
caps missing $p$, each of size at most fixed $b$; $0<\varepsilon\le1$.}
\KwOut{A supplied-query tree of expected cost at most
$(1+\varepsilon)\OPT_E$.}
$N\gets|P|$;
$\ell\gets\max\{2b,b(b-1),\lceil8b^2/\varepsilon\rceil\}$\;
\If{$N<\ell$}{\KwRet $\ExactTree(P\cup\{p\},\mathrm E)$ with raw unit weights\;}
Partition the convex cyclic order of $P$ into pieces of length
$\ell$; mark their first $\min\{b,\text{piece size}\}$ points as $D$\;
$U\gets P$; start an empty miss-spine prefix\;
\While{$D\cap U\ne\varnothing$}{
 Choose $x\in D\cap U$ and a supplied cap $q$ containing $x$\;
 $A\gets U\cap q$\;
 Append $q$ to the prefix; attach a legal local code for $A$
 of height at most $|A|-1$ on its hit side\;
 $U\gets U\setminus A$; continue on the miss side\;
}
Intersect the unmarked pieces with $U$ to obtain the residual blocks\;
Store each physical block's map from its masks to realizing supplied
queries; group identical mask families as types $T$ and compute every
local cost $g_T(A)$ by \ExactTree\;
Compute the initial count vector $c$ of the full residual block masks\;
Evaluate \eqref{eq:uniform-cap-count-dp}, memoizing minimizing transitions\;
Backtrack from $c$: select a physical block in the chosen state,
query that block's stored realizing line, and translate the local
positive code using the same block's query map\;
Attach the reconstructed continuation to the prefix's miss leaf\;
\KwRet the resulting tree\;
\end{algorithm}

The proof uses only the cyclic-interval property of the caps. In
particular it also applies when the nonwitness points alone are in
convex position. For fixed $b$, a constant-factor hardness reduction
through uniform one-layer shallow-cap instances would contradict
this PTAS unless $\mathrm P=\mathrm{NP}$. A construction based on
bounded-size caps must therefore prove interactions beyond this
cyclic-interval subclass to establish such a gap.
